\def\checkmark{\tikz\fill[scale=0.4](0,.35) -- (.25,0) -- (1,.7) -- (.25,.15) -- cycle;}


\newcount\contourmarkcount
\newdimen\contourraise

{\catcode`\|=13
\gdef\installbarmark#1\ignorespaces{%
    #1\ignorespaces%
    \catcode`\|=13%
    \global\contourmarkcount=0\relax%
    \global\def\lastmarkshift{0}%
    \let|=\marktext}%
}


\def\star{*}
\newcommand\marktext[1][*]{%
    \def\markshift{#1}%
    % When not followed by the optional argument
    % the contour mark is set at the previous height.
    \ifx\markshift\star
        \let\markshift=\lastmarkshift%
    \fi
    \global\advance\contourmarkcount by1\relax%
    \xdef\tmpmark{\the\contourmarkcount}%
    \tikz[remember picture, overlay, y=\pgfkeysvalueof{/tikz/contour scale}]
        \path [yshift=\contourraise, shift={(0,\markshift)}]
                coordinate (\contourmarkprefix-\tmpmark);%
    \global\contourmarkcount=\tmpmark\relax% 
    \global\let\lastmarkshift=\markshift%
}


\tikzset{
    intonation contour/.style={%
        execute at begin node={%
            \installbarmark%
        },
        append after command={%
            \pgfextra{%
                \ifnum\contourmarkcount>1
                    \draw [contour] (\contourmarkprefix-1)
                        \foreach \y in {2,...,\the\contourmarkcount}{ -- (\contourmarkprefix-\y) };
                \fi
            }
        }
    },
    % How far above the base line of the text,
    raise contour/.code=\pgfmathsetlength\contourraise{#1},
    % The `scale' for the values in the contour height specification
    contour scale/.initial=3pt,
    % The prefix for the contour marks.
    contour mark prefix/.code=\xdef\contourmarkprefix{#1},
    contour mark prefix=intonation contour,
    contour/.style={
        draw, 
        rounded corners=1ex,
    }           
}

\newcommand{\textelp}{(...)}

\renewcommand{\textlongs}{ſ}

\newcommand{\enquote}[1]{\glqq#1\grqq}


